\subsection{Role-Based Query Scoping}

A natural-language interface that touches a hospital database raises an obvious authorisation question: a doctor must not be able to ask \emph{``how many patients had diabetes last year''} and silently get aggregated counts that include other doctors' panels. The agent solves this by injecting a \texttt{DOCTOR\_SCOPE\_CLAUSE} constraint into both the generation and the validation prompts whenever the authenticated user is a doctor. The clause requires every generated query to restrict results to encounters with the doctor's \texttt{provider\_id}, and, when the question is about patients, to restrict to patients having at least one encounter with that \texttt{provider\_id} (via a documented subquery template). It also forbids exposing other doctors' patients, encounters or aggregate counts, and instructs the agent to refuse the question if it cannot be answered under the scope.

The clause is inserted in the \texttt{scope\_clause} placeholder of both \texttt{PROMPT\_GENERATE\_SQL} and \texttt{PROMPT\_VALIDATE\_SQL}, so the auditor sees the same constraint as the generator. As a last safety net, \texttt{run\_agent(\ldots)} inspects the final SQL one last time: if \texttt{provider\_id} does not appear in the query string, the answer is replaced with a refusal regardless of what the validator concluded.

The two other roles are handled differently. A \texttt{MANAGER} has system-wide visibility and receives no extra clause. A \texttt{PATIENT} never reaches the \texttt{/ask} endpoint at all: the route is closed by the role gate of Section~5.7.1.
